\section{Doctrine, Feasts, and Images}

A biography owes the reader an accurate statement of what the Catholic
Church teaches about Peter's office, kept distinct from the first-century
evidence already presented. The controlling text is the First Vatican
Council's dogmatic constitution \work{Pastor aeternus} (18 July 1870),
checked for this study in its official Latin. Its first chapter teaches
that a primacy of true jurisdiction over the whole Church was promised and
conferred on Peter alone (\textit{beatum Petrum ceteris Apostolis
praeponens}), citing exactly the two texts this biography has quoted:
Matthew 16:16--19 and the risen Christ's ``\textit{Pasce agnos meos: Pasce
oves meas}'' of John 21:15--17. Its second chapter defines that this
primacy endures by Christ's own institution in perpetual successors
(\textit{perpetuos successores}), and that the Roman Pontiff is Peter's
successor in it. The third defines that primacy as ordinary and immediate
jurisdiction over all churches, and the fourth defines, under stated
conditions, the infallibility of the Roman Pontiff's \textit{ex cathedra}
definitions of faith and morals---exercising, the text says, the
infallibility ``with which the divine Redeemer willed his Church to be
endowed,'' promised ``in blessed Peter'' (project paraphrase of the Latin
definition). These are dogmatic determinations about the constitution of
the Church, resting on Scripture read within tradition; the council does
not supply, and does not claim to supply, new historical information about
the first century, and this study reports its teaching as the Church's
authoritative reception of the Petrine texts---the same texts whose
critical history was noted above---without either importing the developed
papal office back into Galilee or subtracting from what the council
actually defined.

Liturgy fixed Peter in the calendar early and permanently. Rome kept the
joint feast of Peter and Paul on 29 June from at least 258, when a
commemoration of the two apostles \textit{ad Catacumbas} is recorded under
that date; the day remains the solemnity of Saints Peter and Paul in the
present General Roman Calendar, with the feast of the Chair of Saint Peter
on 22 February (anciently doubled by a January chair-feast) and the
dedication of the basilicas of Peter and Paul on 18 November. The feasts
commemorate the cult's history, not documented anniversaries. Iconography
settled just as firmly: the keys of Matthew 16:19 became Peter's universal
attribute (and the emblem of the papacy and its flag), joined by the
inverted cross of the martyrdom legend, the cock of the denials, boat and
net and fish; from the fourth century artists gave him the same broad
face, curled hair, and short beard, a convention of recognition, not a
portrait. He is patron of fishermen, of penitents---the tradition loved
the weeping Peter---and of the popes who count from him; churches of
St.~Peter's chains, of the \textit{Quo vadis}, and of the Vatican itself
map his legend onto Rome. In the divided Christian memory he remains
common ground and contested ground at once: all confessions read the same
texts, and the interpretation of the rock and the keys is still the
sharpest single line between them---a question of theology, not of
biography, and left where it belongs.
